\exercisegroup

\begin{enumerate}
\def\labelenumi{\arabic{enumi}.}
\item
  Let I be a directed set, let \((\mathbf{J}_{\lambda})_{\lambda \in \mathbf{L}}\) be a family of subsets of I, indexed by a directed set L, such that (i) for each \(\lambda \in \mathbf{L}\) , \(\mathbf{J}_{\lambda}\) is directed with respect to the induced ordering; (ii) the relation \(\lambda \leqslant \mu\) implies \(\mathbf{J}_{\lambda} \subset \mathbf{J}_{\mu}\) ; (iii) I is the union of the family \((\mathbf{J}_{\lambda})\) . Let \((\mathbf{E}_{\alpha}, f_{\alpha\beta})\) be an inverse system of sets relative to I, let E be its inverse limit, and for each \(\lambda \in \mathbf{L}\) let \(\mathbf{F}_{\lambda}\) be the inverse limit of the system obtained from \((\mathbf{E}_{\alpha}, f_{\alpha\beta})\) by restricting the index set to \(\mathbf{J}_{\lambda}\) . For \(\lambda \leqslant \mu\) let \(g_{\lambda \mu}\) be the canonical mapping of \(\mathbf{F}_{\mu}\) into \(\mathbf{F}_{\lambda}\) (no. 1). Show that \((\mathbf{F}_{\lambda}, g_{\lambda \mu})\) is an inverse system of sets relative to L, and define a canonical bijection of \(\mathbf{F} = \lim \mathbf{F}_{\lambda}\) onto E.
\item
  Let \((\mathbf{E}_{\alpha}, f_{\alpha\beta})\) be an inverse system of sets relative to a directed index set, let \(E = \lim E_{\alpha}\) , and let \(f_{\alpha}: E \to E_{\alpha}\) be the canonical mapping for each \(\alpha\) . Show that, if all the \(f_{\alpha\beta}\) are injective, then \(f_{\alpha}\) is injective.
\item
  Let \((\mathbf{E}_{\alpha}, f_{\alpha\beta})\) and \((\mathbf{F}_{\alpha}, g_{\alpha\beta})\) be two inverse systems of sets relative to the same index set I. For each \(\alpha \in I\) , let \(u_{\alpha}\) be a mapping of \(E_{\alpha}\) into \(F_{\alpha}\) , such that the \(u_{\alpha}\) form an inverse system of mappings. Let \(G_{\alpha} \subset E_{\alpha} \times F_{\alpha}\) be the graph of \(u_{\alpha}\) . Show that \((\mathbf{G}_{\alpha})\) is an inverse system of subsets of \(E_{\alpha} \times F_{\alpha}\) and that its inverse limit may be canonically identified with the graph of \(u = \lim u_{\alpha}\) .
\item
  Let I be a non-empty directed set with no greatest element, and let F be the set of all sequences \(x = (\alpha_1, \alpha_2, \ldots, \alpha_{2n-1}, \alpha_{2n})\) of an even number \(\geqslant 2\) of elements of I with the following properties: (i) \(\alpha_{2i-1} < \alpha_{2i}\) for \(1 \leqslant i \leqslant n\) ; (ii) \(\alpha_{2i-1} \leqslant \alpha_{2j-1}\) for \(1 \leqslant j < i \leqslant n\) . The set F is not empty. Put \(r(x) = \alpha_{2n-1}\) , \(s(x) = \alpha_{2n}\) . The integer \(n\) is called the length of \(x\) .
\end{enumerate}

\begin{enumerate}
\def\labelenumi{(\alph{enumi})}
\tightlist
\item
  For each \(\alpha \in I\) , let \(E_{\alpha}\) be the set of all \(x \in F\) such that \(r(x) = \alpha\) . Then \(E_{\alpha}\) is not empty. For \(\alpha \leqslant \beta\) in \(I\) , we define a mapping \(f_{\alpha \beta}\) of \(E_{\beta}\) into the set of all finite sequences of elements of \(I\) , as follows: if
\end{enumerate}

\[
x = \left(\alpha_ {1}, \dots , \alpha_ {2 n - 1}, \alpha_ {2 n}\right) \in \mathrm{E} _ {\beta},
\]

let \(j\) be the least index such that \(\alpha \leqslant \alpha_{2j - 1}\) ; then

\[
f _ {\alpha \beta} (x) = (\alpha_ {1}, \dots , \alpha_ {2 j - 2}, \alpha , \alpha_ {2 j}).
\]

Show that \(f_{\alpha \beta}(\mathrm{E}_{\beta}) = \mathrm{E}_{\alpha}\) and that \((\mathrm{E}_{\alpha}, f_{\alpha \beta})\) is an inverse system of sets relative to \(\mathbf{I}\) .

\begin{enumerate}
\def\labelenumi{(\alph{enumi})}
\setcounter{enumi}{1}
\item
  Show that if \(x_{\alpha} \in \mathbf{E}_{\alpha}\) and \(x_{\beta} \in \mathbf{E}_{\beta}\) are such that there exists an index \(\gamma\) for which \(\gamma \geqslant \alpha\) and \(\gamma \geqslant \beta\) , and an element \(x_{\gamma} \in \mathbf{E}_{\gamma}\) for which \(x_{\alpha} = f_{\alpha \gamma}(x_{\gamma})\) and \(x_{\beta} = f_{\beta \gamma}(x_{\gamma})\) , then, provided also \(x_{\alpha}\) and \(x_{\beta}\) have the same length, we have \(s(x_{\alpha}) = s(x_{\beta})\) .
\item
  Deduce from (b) that, if \(\mathbf{E} = \lim \mathbf{E}_{\alpha}\) is not empty and if \(y = (x_{\alpha}) \in \mathbf{E}\) , then the set of elements \(s(x_{\alpha})\) is countable and cofinal in \(\mathbf{I}\) .
\item
  Let I be the set of all finite subsets of an uncountable set A, ordered by inclusion. Show that I has no countable cofinal subset, and hence deduce from (c) an example of an inverse system of sets \((\mathbf{E}_{\alpha}, f_{\alpha\beta})\) in which the \(E_{\alpha}\) are non-empty and the \(f_{\alpha\beta}\) are surjective, but for which \(E = \lim E_{\alpha} = \phi\) .
\item
  Deduce from (d) an example of an inverse system of mappings \(u_{\alpha}: \mathbf{E}_{\alpha} \to \mathbf{E}_{\alpha}'\) such that each \(u_{\alpha}\) is surjective but \(\lim u_{\alpha}\) is not surjective (let each \(\mathbf{E}_{\alpha}'\) consist of a single element).
\end{enumerate}

\begin{enumerate}
\def\labelenumi{\arabic{enumi}.}
\setcounter{enumi}{4}
\tightlist
\item
  Let I be a directed set and let \((\mathbf{E}_{\alpha})_{\alpha \in \mathbf{I}}\) be a family of lattices such that each \(\mathbf{E}_{\alpha}\) , endowed with the opposite ordering, is Noetherian (§ 6, no. 5). For each pair \((\alpha, \beta)\) of indices in I such that \(\alpha \leqslant \beta\) let
\end{enumerate}

\[
f _ {\alpha \beta}: \mathbf {E} _ {\beta} \rightarrow \mathbf {E} _ {\alpha}
\]

be an increasing mapping, and suppose that \((\mathbf{E}_{\alpha}, f_{\alpha\beta})\) is an inverse system of sets relative to I. For each \(\alpha \in I\) let \(G_{\alpha}\) be a non-empty subset of \(E_{\alpha}\) such that (i) no two distinct elements of \(G_{\alpha}\) are comparable, (ii) \(f_{\alpha\beta}(G_{\beta}) = G_{\alpha}\) whenever \(\alpha \leqslant \beta\) , (iii) for each \(\alpha \leqslant \beta\) and each \(x_{\alpha} \in G_{\alpha}\) , \(f_{\alpha\beta}(x_{\alpha})\) has a greatest element \(M_{\alpha\beta}(x_{\alpha})\) in \(E_{\beta}\) , (iv) whenever \(\alpha \leqslant \beta\) , if \(h_{\beta} \in E_{\beta}\) is such that there exists \(y_{\beta} \in G_{\beta}\) such that \(y_{\beta} \leqslant h_{\beta}\) , then for each \(x_{\alpha} \in G_{\alpha}\) such that \(x_{\alpha} \leqslant f_{\alpha\beta}(h_{\beta})\) there exists \(x_{\beta} \in G_{\beta}\) such that \(x_{\beta} \leqslant h_{\beta}\) and \(x_{\alpha} = f_{\alpha\beta}(x_{\beta})\) . Under these conditions the inverse limit of the inverse system of subsets \((\mathbf{G}_{\alpha})\) is not empty. The proof runs as follows:

\begin{enumerate}
\def\labelenumi{(\alph{enumi})}
\item
  Let \(J\) be a finite subset of \(I\) . A family \((x_{\alpha})_{\alpha \in J}\) , where \(x_{\alpha} \in G_{\alpha}\) for all \(\alpha \in J\) , is said to be coherent if it satisfies the following two conditions:
\item
  if \(\alpha \in J\) , \(\beta \in J\) , \(\alpha \leqslant \beta\) , then \(x_{\alpha} = f_{\alpha \beta}(x_{\beta})\) ; (ii) for each upper bound \(\gamma\) of \(J\) in \(I\) there exists \(x_{\gamma} \in G_{\gamma}\) such that \(x_{\alpha} = f_{\alpha \gamma}(x_{\gamma})\) for all \(\alpha \in J\) . Show that, for each upper bound \(\gamma\) of \(J\) in \(I\) , the set \(\bigcap_{\alpha \in J}^{-1} f_{\alpha \gamma}(x_{\alpha})\) has a greatest element equal to \(\inf_{\alpha \in J} (M_{\alpha \gamma}(x_{\alpha}))\) ; furthermore, the intersection of \(G_{\gamma}\) and \(\bigcap_{\alpha \in J}^{-1} f_{\alpha \gamma}(x_{\alpha})\) is the set (non-empty by hypothesis) of all \(y_{\gamma} \in G_{\gamma}\) such that
\end{enumerate}

\[
y _ {\gamma} \leqslant \inf _ {\alpha \in J} \left(\mathbf {M} _ {\alpha \gamma} (x _ {\alpha})\right)
\]

(use condition (i)).

\begin{enumerate}
\def\labelenumi{(\alph{enumi})}
\setcounter{enumi}{1}
\tightlist
\item
  Let \(J\) be any subset of \(I\) . A family \(x_{J} = (x_{\alpha})_{\alpha \in J}\) , where \(x_{\alpha} \in G_{\alpha}\) for all \(\alpha \in J\) , is said to be coherent if every finite subfamily of \(x_{J}\) is coherent. If \(J \neq I\) and if \(\beta \in I - J\) , show that there exists \(x_{\beta} \in G_{\beta}\) such that the family \(x_{J \cup \{\beta\}} = (x_{\alpha})_{\alpha \in J \cup \{\beta\}}\) is coherent. (Using (a) and condition (iv), show that, for every finite subset \(F\) of \(J\) , if \(\gamma\) is an upper bound of \(F \cup \{\beta\}\) , then \(f_{\beta \gamma} \left( G_{\gamma} \cap \bigcap_{\alpha \in F} f_{\alpha \gamma}(x_{\alpha}) \right)\) is the (non-empty) set of all \(y_{\beta} \in G_{\beta}\) which are \(\leqslant f_{\beta \gamma} \left( \inf_{\alpha \in F} (M_{\alpha \gamma}(x_{\alpha})) \right)\) . Using the fact that \(E_{\beta}\) endowed with the opposite ordering is Noetherian, show next that there exist a finite subset \(F_0\) of \(J\) and an upper bound \(\gamma_0\) of \(F_0 \cup \{\beta\}\) such that for each finite subset \(F\) of \(J\) and each upper bound \(\gamma\) of \(F \cup \{\beta\}\) we have
\end{enumerate}

\[
f _ {\beta \gamma} \left( \right.\inf _ {\alpha \in \mathbf {F}} \left(\mathrm{M} _ {\alpha j} (x _ {\alpha})\right) \geqslant f _ {\beta \gamma_ {0}} \left(\inf _ {\alpha \in \mathbf {F} _ {0}} \left(\mathrm{M} _ {\alpha \gamma_ {0}} (x _ {\alpha})\right)\right).
\]

Prove then that every element \(x_{\beta} \in \mathbf{F}_{\beta}\) which is \(\leqslant f_{\beta \gamma_0} \left( \inf_{\alpha \in \mathbf{F}_0} (\mathbf{M}_{\alpha \gamma_0}(x_\alpha)) \right)\) satisfies the required conditions.

\begin{enumerate}
\def\labelenumi{(\alph{enumi})}
\setcounter{enumi}{2}
\tightlist
\item
  Finally, complete the proof by showing that there exists a coherent family whose index set is the whole of I. (Order the set of coherent families \(x_{\mathbf{J}}\) by the relation `` \(x_{\mathbf{J}}\) is a subfamily of \(x_{\mathbf{K}}\) '', and apply (b) and Zorn's lemma.)
\end{enumerate}

\begin{enumerate}
\def\labelenumi{\arabic{enumi}.}
\setcounter{enumi}{5}
\item
  Let I be a directed set, and let \((\mathbf{J}_{\lambda})_{\lambda \in \mathbf{L}}\) be a family of subsets of I satisfying the conditions of Exercise 1. Let \((\mathbf{E}_{\alpha}, f_{\beta \alpha})\) be a direct system of sets indexed by I, let \(\mathbf{E} = \lim \mathbf{E}_{\alpha}\) , and for each \(\lambda \in \mathbf{L}\) let \(\mathbf{F}_{\lambda}\) be the direct limit of the direct system obtained from \((\mathbf{E}_{\alpha}, f_{\beta \alpha})\) by restricting the index set to \(J_{\lambda}\) . Whenever \(\lambda \leqslant \mu\) , let \(g_{\mu \lambda}\) be the canonical mapping of \(\mathbf{F}_{\lambda}\) into \(\mathbf{F}_{\mu}\) (no. 6). Show that \((\mathbf{F}_{\lambda}, g_{\mu \lambda})\) is a direct system of sets relative to L, and define a canonical bijection of E onto \(\mathbf{F} = \lim \mathbf{F}_{\lambda}\) .
\item
  Let I be a directed set and let \((\mathbf{E}_{\alpha}, f_{\beta\alpha})\) be a direct system of sets relative to I. For each \(\alpha \in I\) , let \(f_{\alpha}: E_{\alpha} \to E = \lim E_{\alpha}\) be the canonical mapping. In each \(E_{\alpha}\) , let \(R_{\alpha}\) be the equivalence relation \(f_{\alpha}(x) = f_{\alpha}(y)\) .
\end{enumerate}

Show that, whenever \(\alpha \leqslant \beta\) , the mapping \(f_{\theta \alpha}\) is compatible with the equivalence relations \(\mathbf{R}_{\alpha}\) and \(\mathbf{R}_{\beta}\) . Let \(\mathbf{E}_{\alpha}^{\prime} = \mathbf{E}_{\alpha} / \mathbf{R}_{\alpha}\) , and let \(f_{\beta \alpha}^{\prime}\) be the mapping of \(\mathbf{E}_{\alpha}^{\prime}\) into \(\mathbf{E}_{\beta}^{\prime}\) induced by \(f_{\beta \alpha}\) on passing to the quotients. Show that \(f_{\beta \alpha}^{\prime}\) is injective and that \((\mathbf{E}_{\alpha}^{\prime}, f_{\beta \alpha}^{\prime})\) is a direct system of sets, and define a canonical bijection of \(\mathbf{E}\) onto \(\lim_{\alpha \to 0} \mathbf{E}_{\alpha}^{\prime}\) .

\begin{enumerate}
\def\labelenumi{\arabic{enumi}.}
\setcounter{enumi}{7}
\item
  Let \((\mathbf{E}_{\alpha}, f_{\beta\alpha})\) and \((\mathbf{F}_{\alpha}, g_{\beta\alpha})\) be two direct systems of sets, both indexed by the same directed set I. For each \(\alpha \in I\) , let \(u_{\alpha}\) be a mapping of \(E_{\alpha}\) into \(F_{\alpha}\) , such that the \(u_{\alpha}\) form a direct system of mappings. Let \(G_{\alpha} \subset E_{\alpha} \times F_{\alpha}\) be the graph of \(u_{\alpha}\) . Show that \((\mathbf{G}_{\alpha})\) is a direct system of subsets of \(E_{\alpha} \times F_{\alpha}\) and that its direct limit may be canonically identified with the graph of \(u = \lim u_{\alpha}\) .
\item
  Let I be an arbitrary preordered set, and let \((\mathrm{E}_{\alpha})_{\alpha\in\mathbf{I}}\) be a family of sets indexed by I. For each pair of indices \((\alpha,\beta)\) such that \(\alpha\leqslant\beta\) , let \(f_{\beta\alpha}\) be a mapping of \(E_{\alpha}\) into \(E_{\beta}\) , and suppose that these mappings satisfy conditions \((\mathrm{LI}_{\mathbf{I}})\) and \((\mathrm{LI}_{\mathbf{II}})\) . Let G be the set which is the sum of the family \((\mathrm{E}_{\alpha})\) and (with the notation of no. 5) let \(R\{x,y\}\) be the relation `` \(\lambda(x)=\alpha\leqslant\lambda(y)=\beta\) and \(y=f_{\beta\alpha}(x)\) '' between two elements x, y of G. Let \(R'\) be the equivalence relation on G whose graph is the smallest of the graphs of equivalence relations which contain the graph of R (Chapter II, §6, Exercise 10). The set \(E=G/R'\) is called the direct limit of the family \((\mathrm{E}_{\alpha})\) with respect to the family of mappings \((f_{\beta\alpha})\) , and we write \(E=\lim_{\rightarrow}E_{\alpha}\) . When the index set I is directed, show that this definition agrees with that given in no. 5. In the general case, the restriction to \(E_{\alpha}\) of the canonical mapping of G into E is called the canonical mapping of \(E_{\alpha}\) into E and is denoted by \(f_{\alpha}\) . Suppose we are given, for each \(\alpha\in I\) , a mapping \(u_{\alpha}\) of \(E_{\alpha}\) into F such that \(u_{\beta}\circ f_{\beta\alpha}=u_{\alpha}\) whenever \(\alpha\leqslant\beta\) ; show that there exists a unique mapping u of E into F such that \(u=u_{\alpha}\circ f_{\alpha}\) for each \(\alpha\in I\) .
\end{enumerate}
